% ══════════════════════════════════════════════════════════════
%  SECTION 2: INSTITUTIONAL BACKGROUND AND IDENTIFICATION
%  Draft: v3_institutional.tex
%  Target: paper/v2/main.tex Section 2
% ══════════════════════════════════════════════════════════════

\section{Institutional Background and Identification}
\label{sec:background}

\subsection{The Policy: An Enforcement Shift, Not a New Rule}

Association football (soccer) matches are governed by seventeen Laws of the
Game, maintained by the International Football Association Board (IFAB).
Law~7 specifies the duration of a match: two periods of 45~minutes each,
with ``allowance'' made ``in each half for all time lost'' due to
substitutions, injuries, time-wasting, disciplinary sanctions, VAR (video
assistant referee) reviews, goal celebrations, and ``any other cause''
\citep{ifab2023laws}. Unlike American football or basketball, the match
clock runs continuously; time lost to stoppages is restored in a single
block of ``added time'' (also called ``stoppage time'') appended to the end
of each half. The referee signals the minimum number of added minutes; play
continues until the referee ends the half at discretion.

The critical institutional fact is that Law~7 has \textit{always} required
time restoration. The rule did not change in 2022; what changed was
enforcement. For decades, referees applied a convention---typically one to
four minutes of added time per half---that systematically undercompensated
for actual time lost, sustained by a norm of approximation rather than
precise accounting.

At the 2022 FIFA World Cup in Qatar (November--December 2022), Pierluigi
Collina, chairman of FIFA's Referees Committee, instructed match officials
to add the \textit{actual} time lost rather than the conventional
approximation. Several group-stage matches exceeded 100~minutes of total
playing time---England--Iran ran to 117~minutes---generating worldwide
media attention \citep{wei2024goal, lames2024goal}. IFAB endorsed the
enforcement expectation at its March 2023 Annual General Meeting, stating
that referees ``should'' restore the full time lost. The hortatory language
(``should'' not ``must'') granted national federations discretion. UEFA
explicitly opted out for the Champions League and Europa League.

This distinction---between a new rule and a new enforcement norm for an
existing rule---matters for interpretation. The treatment we study is not
the introduction of a novel policy instrument but the activation of a
dormant one. The closest analogy in the regulation literature is an
enforcement crackdown rather than new legislation \citep{johnson2020compliance,
greenstone2004did}. The statutory authority was already in place; what
changed was the expectation of compliance.

\subsection{League Adoption: A Simultaneous Shock with Heterogeneous Local
Enforcement}

The IFAB guidance was binding on all 211 FIFA member associations from the
2023--24 season. In practice, enforcement intensity varied dramatically
across leagues---and, crucially, \textit{within} countries. Our sample
covers 16 European leagues over 12 seasons (2014--15 through 2025--26),
encompassing 36,267 matches.

Table~\ref{tab:enforcement} reports the league sample. Five leagues---the
Premier League (England, E0), La~Liga (Spain, SP1), Ligue~1 (France, F1),
the Bundesliga (Germany, D1), and Serie~A (Italy, I1)---are classified as
treated. All five are top divisions of the ``Big Five'' European football
markets, sharing direct institutional ties to FIFA, intense media scrutiny,
and national refereeing organizations that issued explicit enforcement
guidance.

We are transparent about adoption timing. The Premier League, La~Liga,
and Ligue~1 opened their 2023--24 seasons on August~11, 2023; the
Bundesliga on August~18; Serie~A on August~19. All five adopted within an
eight-day window. This is a \textit{simultaneous} treatment shock, not a
staggered rollout; we do not claim otherwise. The eight-day spread is
available as a minor robustness check but provides no meaningful
identifying variation.

The remaining eleven leagues serve as controls:

\begin{itemize}[leftmargin=*]
  \item \textit{Second divisions of treated countries} (5):
    Championship (E1), 2.~Bundesliga (D2), Ligue~2 (F2), Serie~B (I2),
    and Segunda~Divisi\'{o}n (SP2). These share a federation, calendar,
    and---to varying degrees---a referee pool with their top-flight
    counterparts, but did not receive the same enforcement emphasis.
  \item \textit{Other top flights} (6): Eredivisie (N1), Primeira~Liga
    (P1), S\"{u}per~Lig (T1), Belgian Pro~League (B1), Greek
    Super~League (G1), and Scottish Premiership (SC0).
\end{itemize}

\begin{table}[t]
\centering
\caption{Enforcement Intensity by League}
\label{tab:enforcement}
\small
\begin{threeparttable}
\begin{tabular}{llccccl}
\toprule
League & Code & Pre Mean & Post Mean & $\Delta$ & Cohen's $d$ & Intensity \\
\midrule
Premier League$^\dagger$ & E0 & 5.39 & 7.50 & +2.11 & 1.13 & High \\
Ligue 1$^\dagger$ & F1 & 4.71 & 6.50 & +1.79 & 1.05 & High \\
Bundesliga$^\dagger$ & D1 & 4.34 & 6.12 & +1.78 & 0.92 & High \\
Scottish Premiership & SC0 & 2.83 & 4.30 & +1.47 & 0.83 & High \\
La Liga$^\dagger$ & SP1 & 5.18 & 6.64 & +1.46 & 0.70 & High \\
Belgian Pro League & B1 & 3.09 & 4.40 & +1.31 & 0.66 & High \\
2. Bundesliga & D2 & 3.95 & 4.98 & +1.04 & 0.59 & High \\
Championship & E1 & 5.67 & 6.70 & +1.03 & 0.59 & High \\
Primeira Liga & P1 & 5.97 & 6.94 & +0.97 & 0.42 & Moderate \\
Ligue 2 & F2 & 2.81 & 3.76 & +0.95 & 0.56 & Moderate \\
Greek Super League & G1 & 3.58 & 4.41 & +0.82 & 0.38 & Moderate \\
Serie A$^\dagger$ & I1 & 5.18 & 5.91 & +0.73 & 0.40 & Moderate \\
Eredivisie & N1 & 5.24 & 5.71 & +0.47 & 0.23 & Low \\
Serie B & I2 & 4.05 & 3.97 & -0.08 & -0.04 & Low \\
\bottomrule
\end{tabular}
\begin{tablenotes}\small
\item \textit{Notes:} Pre = seasons before 2023-24. Post = 2023-24 onward. $\Delta$ = Post $-$ Pre mean second-half stoppage time (minutes). Cohen's $d$ = standardized effect size. Intensity classified as High ($\Delta > 1.0$), Moderate ($0.5 < \Delta \leq 1.0$), Low ($\Delta \leq 0.5$). $^\dagger$ denotes Big Five league.
\end{tablenotes}
\end{threeparttable}
\end{table}


The second-division controls are a particular strength. If the Premier
League increased its stoppage time but the Championship---in the same
country, same season, same football culture---did not, the within-country
contrast isolates the directive's enforcement channel from country-level
confounds.

Pre-directive, treated leagues averaged 4.99~minutes of second-half
stoppage ($N = 10{,}763$); controls averaged 4.32~minutes ($N = 15{,}865$).
Post-directive, treated leagues rose to 6.56~minutes ($+1.57$), controls to
5.62~minutes ($+1.30$)---the control-group increase reflecting both secular
trends and partial directive transmission through shared refereeing bodies.

\subsection{Identification Strategy}
\label{subsec:identification}

Partial treatment of the control group is the central identification
challenge. Several control leagues---notably the Championship and
2.~Bundesliga, which share referee pools with their top-flight
counterparts---also increased stoppage time. A naive comparison therefore
understates the true effect. We address this with three complementary
strategies.

\paragraph{Primary: Treated versus never-treated DiD.}
We pool all 16 leagues and estimate the differential increase via a
two-period DiD with league and season fixed effects, clustering at the
league level with wild cluster bootstrap inference
\citep{cameron2008bootstrap}. Because some control leagues partially
adopted the directive, this estimate is a \textit{conservative lower
bound}. The Callaway--Sant'Anna estimator \citep{callaway2021difference}
is appropriate given binary treatment timing; we verify robustness to
\citet{sunAbraham2021} and \citet{borusyak2024revisiting}.

\paragraph{Co-primary: Within-country pairs.}
We restrict the sample to four countries with both a treated top-flight
and a control second division: England (E0/E1), Germany (D1/D2), France
(F1/F2), and Italy (I1/I2). Spain is excluded because Segunda~Divisi\'{o}n
data are unavailable post-2022--23. The estimating equation is:
\begin{equation}\label{eq:within}
  Y_{ilct} = \alpha + \beta \cdot (\text{Post}_t \times \text{TopFlight}_l)
  + \gamma_{ct} + \mu_l + \varepsilon_{ilct},
\end{equation}
where $\gamma_{ct}$ denotes country$\times$season fixed effects and
$\mu_l$ league fixed effects. These pairs share a federation, calendar,
broadcasting environment, and---crucially---partially overlapping referee
labor markets. The design absorbs country-level trends, isolating
differential enforcement transmission. With only 8 league clusters,
precision is reduced but identification is tighter.

\paragraph{Supporting: Referee-level dose variation.}
Within the treated leagues, 407 referees officiated matches in our sample.
Among the 85 with at least 20 matches in both pre- and post-directive
periods, 100\% increased their mean stoppage time. Cross-referee variation
in pre-directive baselines provides continuous dose variation: referees
already generous with stoppage had less room to adjust than those who
previously added minimal time. This within-league strategy does not depend
on the cross-league contrast, though it requires careful treatment of
referee selection and assignment.

\subsection{Pre-Trends and the Limits of the Event Study}
\label{subsec:pretrends}

The parallel trends assumption is the Achilles' heel of this design. The
event study (Section~\ref{subsec:pretrends_results}) rejects the null of
parallel pre-trends ($F = 9.0$, $p < 0.001$). We do not minimize this.
A reader who requires classical parallel-trends validation will not find it
here.

The \citet{rambachan2023more} HonestDiD sensitivity analysis confirms the
fragility: the breakdown value is $\bar{M}^* = 0.11$, meaning even modest
extrapolation of pre-existing differential trends eliminates significance.
For individual post-treatment periods, breakdown values are negative---the
honest confidence intervals include zero under any positive trend
extrapolation.

Why, then, do we maintain a causal interpretation? The answer is
institutional, not statistical.

\textit{First}, the within-country pairs absorb country-level secular
trends. If English stoppage norms were drifting upward for both the
Premier League and the Championship, the country$\times$season fixed
effects in Equation~\eqref{eq:within} net this out. The within-country
estimate ($\hat{\beta} = 0.87$, $p = 0.033$) identifies off the
\textit{differential} enforcement within the same country and season.

\textit{Second}, the mechanism is directly observable. The directive was
publicly announced, its transmission through refereeing organizations is
documented, and referee-level compliance is measurable. Treatment existence
and timing need not be inferred from the data.

\textit{Third}, the pre-trend failure has a plausible explanation. Stoppage
norms evolved heterogeneously across countries during the pre-treatment
period, driven by staggered VAR adoption, changes in match tempo, and
culturally specific attitudes toward time management. These secular trends
generate pre-treatment coefficient drift in the cross-country event study
without necessarily violating the within-country identifying assumption.

We present multiple strategies not to cherry-pick the most favorable but
because no single design is dispositive. The directive produced a discrete
shift in stoppage time that is robust in sign and approximate magnitude
across all specifications, but the precise causal estimate depends on
assumptions about pre-trends the data cannot fully validate. Readers who
require $\bar{M}^* > 0$ should interpret our results as precise descriptive
accounting of what changed when enforcement norms shifted
\citep{roth2022pretest}.
